\section{Conclusiones}

\setlength{\parindent}{1.5cm}

\indent El desarrollo del presente proyecto permitió al equipo comprender de manera rigurosa y
práctica el papel que desempeña la \textbf{colección canónica de conjuntos de ítems LR(0)}
en la construcción de analizadores sintácticos ascendentes. A diferencia del análisis
predictivo descendente LL(1) \textemdash{}estudiado en la asignación anterior mediante los
conjuntos \textsc{Primero} y \textsc{Siguiente}\textemdash{} el análisis LR(0) no necesita
predecir la producción a aplicar: en cambio, acumula información en la pila a medida que
desplaza símbolos de la entrada y decide reducir cuando el estado activo contiene un ítem
completo. La colección canónica es precisamente el catálogo exhaustivo de todos los estados
posibles de ese proceso.

\indent La implementación de los algoritmos de clausura e ir\_a puso de manifiesto la
naturaleza composicional del autómata LR(0): \texttt{closure} expande horizontalmente un
conjunto de ítems capturando todas las predicciones posibles para un no terminal, mientras
que \texttt{goto} desplaza el punto a través de un símbolo y vuelve a clausurar, generando
un nuevo estado. La correcta composición iterativa de ambas funciones en
\texttt{construir\_coleccion} garantiza que la colección sea mínima (sin estados duplicados)
y completa (cubre todos los prefijos viables). El uso de \texttt{frozenset} como
representación canónica de cada estado fue clave para detectar estados ya conocidos en
tiempo constante.

\indent El diseño en tres capas \textemdash{}modelos, persistencia y
presentación\textemdash{} demostró su valor durante la integración: cada sub equipo pudo
desarrollar y verificar su módulo de forma autónoma antes de conectarlo con los demás, ya
que las interfaces entre capas estaban definidas por las \textit{dataclasses} inmutables.
El patrón fachada aplicado en \texttt{algoritmo\_lr0.py} simplificó la capa de presentación
al exponer solo las funciones relevantes y ocultar los detalles internos del módulo
algorítmico.

\indent El sistema de traza detallada incorporado en \texttt{construir\_coleccion} resultó
especialmente valioso en el contexto educativo: al registrar cada llamada a \texttt{closure}
e \texttt{ir\_a} con su resultado, el usuario puede seguir paso a paso la construcción del
autómata y confrontarla con los cálculos manuales realizados en clase. Esta característica
convierte la herramienta en un instrumento de verificación, no solo de cómputo.

\indent Como perspectiva futura, la arquitectura del sistema facilita la extensión hacia
técnicas más potentes: los conjuntos \textsc{Siguiente} \textemdash{}calculados en la
asignación anterior\textemdash{} pueden incorporarse directamente para construir tablas
SLR(1), y con la adición de lookaheads por ítem se puede escalar a LALR(1) o LR(1). La
modularidad ya establecida minimizaría el esfuerzo de esas extensiones.

\indent En síntesis, el proyecto consolidó la comprensión de la fase ascendente del diseño
de compiladores y reforzó buenas prácticas de ingeniería de software \textemdash{}estructuras
inmutables, separación de responsabilidades y pruebas incrementales\textemdash{} que serán
aplicadas en las etapas subsecuentes del curso.
